\chapter{Теоретические сведения}
\label{ch:theory}

\section{Стек LAMP}

\textbf{LAMP} --- аббревиатура от Linux, Apache, MySQL/MariaDB, PHP.
Это классический набор программного обеспечения для запуска динамических
веб-приложений на серверах под управлением Linux.

Компоненты стека:
\begin{itemize}
  \item \textbf{Linux} --- операционная система; в данной работе
        используется Ubuntu~22.04~LTS;
  \item \textbf{Apache2} --- HTTP-сервер; обрабатывает входящие
        запросы и передаёт их PHP;
  \item \textbf{MariaDB} --- реляционная СУБД, fork MySQL;
        хранит данные WordPress (посты, пользователи, настройки);
  \item \textbf{PHP} --- серверный скриптовый язык; WordPress
        написан на PHP и требует версии не ниже~7.4.
\end{itemize}

Установить полный стек одной командой позволяет утилита
\texttt{tasksel}:
\begin{lstlisting}
apt-get install -y tasksel
tasksel install lamp-server
\end{lstlisting}

\section{Apache2}

\textbf{Apache HTTP Server} (\texttt{apache2}) --- свободный
кроссплатформенный веб-сервер. Конфигурация хранится в
\texttt{/etc/apache2/}. Ключевые файлы и каталоги:
\begin{itemize}
  \item \texttt{apache2.conf} --- глобальные настройки;
  \item \texttt{sites-available/} --- файлы виртуальных хостов;
  \item \texttt{sites-enabled/} --- симлинки на активные хосты;
  \item \texttt{mods-enabled/} --- подключённые модули.
\end{itemize}

Каждый \textbf{Virtual Host} описывает отдельный сайт. Пример минимальной
конфигурации виртуального хоста для WordPress:
\begin{lstlisting}
<VirtualHost *:80>
    ServerName 192.168.N.6
    DocumentRoot /var/www/html
    <Directory /var/www/html>
        AllowOverride All
    </Directory>
</VirtualHost>
\end{lstlisting}

Модуль \texttt{rewrite} необходим WordPress для работы
<<красивых>> постоянных ссылок:
\begin{lstlisting}
a2enmod rewrite
systemctl restart apache2
\end{lstlisting}

\section{MariaDB}

\textbf{MariaDB} --- открытая СУБД, полностью совместимая с MySQL.
После установки необходимо:
\begin{enumerate}
  \item Установить пароль root-пользователя:
        \texttt{mysqladmin -u root password 'ПАРОЛЬ'}
  \item Создать базу данных для WordPress:
\begin{lstlisting}
CREATE DATABASE wordpress CHARACTER SET utf8 COLLATE utf8_bin;
\end{lstlisting}
  \item Создать пользователя и выдать привилегии:
\begin{lstlisting}
CREATE USER 'author'@'localhost' IDENTIFIED BY 'Pssw0rd';
GRANT ALL PRIVILEGES ON wordpress.* TO 'author'@'localhost';
FLUSH PRIVILEGES;
\end{lstlisting}
\end{enumerate}

\section{WordPress}

\textbf{WordPress} --- CMS с открытым исходным кодом. Распространяется
архивом \texttt{.tar.gz}. Основные шаги установки:
\begin{enumerate}
  \item Скачать архив с официального сайта:
\begin{lstlisting}
wget https://ru.wordpress.org/latest-ruRU.tar.gz
tar xzvf latest-ruRU.tar.gz
\end{lstlisting}
  \item Скопировать пример конфигурационного файла:
\begin{lstlisting}
cp wp-config-sample.php wp-config.php
\end{lstlisting}
  \item Прописать параметры базы данных в \texttt{wp-config.php}:
\begin{lstlisting}
define('DB_NAME',     'wordpress');
define('DB_USER',     'author');
define('DB_PASSWORD', 'Pssw0rd');
define('DB_HOST',     'localhost');
\end{lstlisting}
  \item Скопировать файлы в корень сайта Apache:
\begin{lstlisting}
rsync -aq wordpress/ /var/www/html/
chown -R www-data:www-data /var/www/html
\end{lstlisting}
  \item Завершить установку интерактивно через браузер,
        перейдя по адресу \texttt{http://IP-СЕРВЕРА/}.
\end{enumerate}

После завершения мастера установки WordPress создаёт необходимые
таблицы в базе данных и становится доступным для входа.

\endinput
